
\sloppy
\chapter*{\centering Abstract}
\onehalfspacing

The rapid spread of fake news on digital media makes manual fact-checking impractical, creating a need for automated detection systems. This thesis proposes a fake news detection framework that combines NLP-based text preprocessing with a comparative evaluation of three classification models: a standalone BiLSTM network, a fine-tuned BERT model, and a hybrid BERT + BiLSTM model.

The models were evaluated on the LIAR and ISOT benchmark datasets. BiLSTM performed poorly, reaching only 50\% accuracy, while BERT achieved 95\% accuracy with strong precision and recall. The hybrid BERT + BiLSTM model matched BERT's performance but did not exceed it, showing that the added sequential layer offered no further benefit in this setup.

These results confirm that transformer-based models clearly outperform sequential deep learning models for fake news detection and are well suited to practical use in social media monitoring and fact-checking. Future work will explore Large Language Models, Retrieval-Augmented Generation, and Explainable AI to further improve accuracy and interpretability.

\vspace{0.5in}
\noindent
{\bf Keywords:} Fake News Detection, Natural Language Processing, Machine Learning, Deep Learning, BERT, BiLSTM, Transformers, Explainable AI, Misinformation.
